% =============================================================================
%                  WEEK 5: MID-LEVEL VISION --- SEGMENTATION
% =============================================================================
\section{Week 5: Mid-Level Vision --- Segmentation}

\subsection{Topic Summary}

\begin{keybox}[Learning Objectives]
By the end of this week, you will be able to:
\begin{itemize}
  \item Define the role of mid-level vision and explain feature space and similarity/distance measures.
  \item Apply thresholding (single, hysteresis, local) and morphological clean-up (erosion, dilation, opening, closing).
  \item Apply region-based segmentation (region growing, region merging, split-and-merge) on small numerical examples.
  \item Apply clustering for segmentation (k-means and agglomerative hierarchical clustering with single/complete/average/centroid linkage).
  \item Explain segmentation by graph cutting and the Normalised Cuts (Ncuts) criterion.
  \item Apply the Hough transform for straight lines and explain its generalisation to circles and arbitrary shapes.
  \item Describe active contours (snakes) and the energy terms they minimise.
\end{itemize}
\end{keybox}

\textbf{Big Picture.}
Mid-level vision is about \term{grouping}: take low-level image elements (pixels, edges, filter responses) and gather those that ``belong together'' into segments. Every method in this week is a different answer to the same question: which feature(s) define belonging, and how do we measure similarity? Each pixel becomes a point in a \term{feature space}, and segmentation reduces to partitioning that space. Methods differ in (i) what feature/measure they use, (ii) whether they reason in feature space alone or use spatial information, and (iii) what model they assume (uniform region, cluster, parametric shape, smooth curve).

\separator

\textbf{What Examiners Like to Ask:}
\begin{itemize}
  \item ``Briefly describe the role of mid-level vision.''
  \item How to choose a threshold from an intensity histogram; explain hysteresis thresholding.
  \item Apply hysteresis thresholding by hand on a small numerical image.
  \item Show the result of erosion-then-dilation (opening) and dilation-then-erosion (closing) for 4- and 8-connected neighbourhoods.
  \item Write pseudo-code for region growing / region merging / split-and-merge / k-means / agglomerative / Ncuts / Hough.
  \item Apply each algorithm to a 3$\times$3 feature-vector image with SAD as similarity and a numerical threshold.
  \item Describe single-link / complete-link / group-average / centroid linkage.
  \item Compute a Hough accumulator array for a small binary image at given $\theta$ values.
  \item Explain the two energy terms used by an active contour (snake).
\end{itemize}

% =============================================================================
\subsection{Core Concepts}

\textbf{Roadmap:}
\begin{enumerate}
  \item Feature space and similarity / distance measures
  \item Thresholding (and morphological clean-up)
  \item Region-based methods (growing, merging, split-and-merge)
  \item Clustering (k-means, agglomerative hierarchical, graph cuts/Ncuts)
  \item Fitting (Hough transform, active contours)
\end{enumerate}

\bigseparator

% -----------------------------------------------------------------------------
\subsubsection{Feature Space and Similarity}

\begin{defbox}[Mid-Level Vision (Segmentation)]
The task of grouping together image elements that ``belong together'' (e.g. lie on the same object) and segmenting them from all other image elements.
\end{defbox}

\textbf{Common features used for grouping} (each motivates a Gestalt grouping law):
\begin{itemize}
  \item Location (proximity), colour/intensity (similarity), texture (similarity)
  \item Depth, motion (common fate)
  \item Not separated by a contour (common region)
  \item Top-down: forms a known shape, or CNN feature vectors
\end{itemize}

\separator

\begin{defbox}[Feature Space]
Each image element is represented as a vector $\mathbf{a} = (a_1, \dots, a_n)$ in an $n$-dimensional feature space. Similarity (and hence grouping) is determined by the distance between points in this space.
\end{defbox}

\textbf{Similarity measures} (high when vectors agree):
\[
\text{Affinity} = \exp\!\left(-\frac{\sum_i (a_i - b_i)^2}{2\sigma^2}\right), \qquad
\text{Cross-correlation} = \sum_i a_i b_i
\]
\[
\text{NCC} = \frac{\sum_i a_i b_i}{\sqrt{\sum_i a_i^2}\,\sqrt{\sum_i b_i^2}}, \qquad
\text{Corr.\ coeff.} = \frac{\sum_i (a_i - \bar a)(b_i - \bar b)}{\sqrt{\sum_i (a_i - \bar a)^2}\,\sqrt{\sum_i (b_i - \bar b)^2}}
\]

\textbf{Distance measures} (low when vectors agree):
\[
\text{Euclidean} = \sqrt{\textstyle\sum_i (a_i - b_i)^2}, \quad
\text{SSD} = \sum_i (a_i - b_i)^2, \quad
\text{SAD} = \sum_i |a_i - b_i|
\]

\begin{tipbox}[Exam Tip: SAD is the default in tutorial questions]
Almost every numerical tutorial question for this week uses \textbf{Sum of Absolute Differences (SAD)} as the similarity measure with a threshold on SAD. Learn to compute SAD quickly: just add the absolute element-wise differences.
\end{tipbox}

\separator

\textbf{Feature scaling and weighting.} Different features have different scales (e.g. pixel coordinates in $[0,2000]$ vs.\ intensity in $[0,1]$). If left raw, large-scale features dominate the distance and small-scale features look ``always similar''. Solutions:
\begin{itemize}
  \item Normalise each feature to a common range.
  \item Use weighted distance, e.g. $\sqrt{w_1(a_1{-}b_1)^2 + w_2(a_2{-}b_2)^2 + \dots}$. Choosing $w_i$ that work well is non-trivial (requires optimisation).
\end{itemize}

\bigseparator

% -----------------------------------------------------------------------------
\subsubsection{Thresholding}

\begin{defbox}[Thresholding]
A 1-D feature space segmentation: assign pixel $(x,y)$ to foreground if its feature value $I(x,y)$ exceeds a threshold:
\[ I'(x,y) = \begin{cases} 1 & \text{if } I(x,y) > t \\ 0 & \text{otherwise} \end{cases} \]
Produces \term{figure-ground} segmentation. The feature can be intensity, colour channel, or output of a filter (e.g. edge magnitude).
\end{defbox}

\textbf{How to choose the threshold:}
\begin{enumerate}
  \item \textbf{Average intensity} of the image.
  \item \textbf{Histogram-based}: plot the intensity histogram, assume it is \term{bimodal}, place the threshold in the valley between the two peaks.
  \item \textbf{Hysteresis thresholding}: use \emph{two} thresholds (one each side of the valley).
  \item \textbf{Fraction-of-image}: choose $t$ so a known fraction of pixels are foreground (when object size is known).
\end{enumerate}

\begin{defbox}[Hysteresis Thresholding]
Define a high threshold $t_h$ and a low threshold $t_l$ (with $t_l < t_h$):
\begin{itemize}
  \item Pixels with $I > t_h$: classified as figure (foreground).
  \item Pixels with $I < t_l$: classified as background.
  \item Pixels with $t_l \le I \le t_h$: classified as figure \emph{only if} they are adjacent to a pixel already classified as figure (otherwise background).
\end{itemize}
\end{defbox}

\textbf{Local / block thresholding.} A single global threshold rarely works for whole images (uneven illumination, shadows). Split the image into rectangular blocks and apply a different threshold per block --- it is the \emph{local} contrast that matters.

\separator

\textbf{Thresholding errors and morphological clean-up.}
Thresholding usually leaves artefacts: missing pixels (holes in objects) and unwanted pixels (specks of noise). Morphological operations clean these up.

\begin{defbox}[Dilation and Erosion]
Both use a \term{structuring element} (typically 3$\times$3) defining a neighbourhood (4-connected: H/V; 8-connected: H/V/diagonal).
\begin{itemize}
  \item \textbf{Dilation}: every background pixel that has a foreground neighbour becomes foreground (1). Expands foreground; fills holes/gaps.
  \item \textbf{Erosion}: every foreground pixel that has a background neighbour becomes background (0). Shrinks foreground; removes bridges/specks.
\end{itemize}
\end{defbox}

\begin{thmbox}[Opening and Closing]
\begin{itemize}
  \item \textbf{Opening} = erosion then dilation. Removes background noise (small foreground blobs disappear); leaves overall shape unchanged.
  \item \textbf{Closing} = dilation then erosion. Fills foreground holes/gaps; leaves overall shape unchanged.
\end{itemize}
\end{thmbox}

\begin{tipbox}[Exam Tip: Opening vs Closing]
Mnemonic: ``open removes specks, close fills holes''. The first operation (erosion or dilation) tells you what is removed; the second restores the surviving objects to roughly their original size.
\end{tipbox}

\bigseparator

% -----------------------------------------------------------------------------
\subsubsection{Region-Based Segmentation}

\textbf{Motivation.} Thresholding ignores spatial information. Region-based methods consider pixel \emph{location} as well as feature values, so the feature space is multi-dimensional (e.g. RGB, or RGB + position).

\separator

\begin{defbox}[Region Growing]
Start from a seed pixel; iteratively add neighbouring pixels whose features are within a similarity threshold of the seed's region.
\end{defbox}

\textbf{Pseudo-code (region growing):}
\begin{enumerate}
  \item While not all pixels have a region label:
  \begin{enumerate}
    \item Choose an unlabelled pixel at random.
    \item Assign it the next unused region label.
    \item For all neighbouring unlabelled pixels:
    \begin{enumerate}
      \item Calculate similarity to the chosen pixel.
      \item If within similarity threshold: give it the same region label, then make that pixel the chosen pixel and recurse.
    \end{enumerate}
  \end{enumerate}
\end{enumerate}

\textbf{When to use.} You have plausible seed locations (or can pick them randomly) and a per-pixel similarity criterion.\\
\textbf{Pitfalls.} Regions can ``leak'' through a single weak spot in the boundary. If similarity uses intensity only, regions stop at edges --- but only if edges are unbroken.

\separator

\begin{defbox}[Region Merging]
Start with every pixel as its own region; iteratively merge adjacent regions whose properties (mean feature vector) are similar.
\end{defbox}

\textbf{Pseudo-code (region merging):}
\begin{enumerate}
  \item Give every pixel a unique region label.
  \item Randomly choose a region not yet marked \emph{final}.
  \begin{enumerate}
    \item Compare its properties with each neighbour's.
    \item For all neighbours within similarity threshold: relabel them to the chosen region's label; recompute properties as the mean of all constituent pixels.
    \item If no neighbour can be merged: mark the chosen region as \emph{final}.
  \end{enumerate}
  \item Repeat until all regions are marked \emph{final}.
\end{enumerate}

\textbf{Pitfall: order-dependence.} The result depends on the order in which regions are merged, because the merged region's mean shifts with each merge.

\separator

\begin{defbox}[Split-and-Merge]
Top-down then bottom-up. Recursively split heterogeneous regions into quadrants (a quadtree) until each region is homogeneous, then merge similar adjacent regions.
\end{defbox}

\textbf{Pseudo-code (split-and-merge):}
\begin{enumerate}
  \item Label every pixel with the same region label.
  \item For each region: if not all pixels are similar, split into 4 quadrants with new labels.
  \item Repeat step 2 until all regions are homogeneous.
  \item For each region: compare to neighbours; merge any whose properties are within the similarity threshold (recompute mean).
  \item Repeat step 4 until no more merges possible.
\end{enumerate}

\textbf{Note.} If image dimensions are not powers of 2, pad the image so it can be quartered cleanly.

\separator

\begin{tipbox}[Exam Tip: General problems with region-based methods]
``Meaningful'' regions are rarely uniform: brightness/colour vary due to lighting, curvature, and texture. Real images are almost never composed of uniform regions of similar intensity/colour. Be ready to state this as a limitation.
\end{tipbox}

\bigseparator

% -----------------------------------------------------------------------------
\subsubsection{Clustering Methods}

\begin{defbox}[Clustering]
A general class of methods for \term{unsupervised classification} of data --- classes are not predefined. Aim: minimise intra-cluster distance, maximise inter-cluster distance.
\end{defbox}

Two sub-classes:
\begin{itemize}
  \item \textbf{Partitional}: divide data into non-overlapping subsets (e.g.\ k-means).
  \item \textbf{Hierarchical}: nested clusters organised as a tree (dendrogram) --- agglomerative (bottom-up) or divisive (top-down).
\end{itemize}

\separator

\begin{defbox}[k-means Clustering]
A partitional algorithm assuming there are $k$ clusters. Iteratively assigns points to the nearest cluster centre, then recomputes each centre as the mean of its assigned points, until centres stop moving.
\end{defbox}

\textbf{Pseudo-code (k-means):}
\begin{enumerate}
  \item Randomly choose $k$ points to act as cluster centres.
  \item For each data point: compute similarity to each cluster centre; allocate to closest.
  \item For each cluster centre: set new position to mean of its allocated points.
  \item Repeat 2--3 until cluster centres are unchanged (converged).
\end{enumerate}

\textbf{Problems with k-means:}
\begin{itemize}
  \item Results depend on initial centre positions (different runs $\Rightarrow$ different clusterings).
  \item Poor when true clusters have very different sizes.
  \item Poor when true clusters are not roughly ``globular''.
  \item Must choose $k$ in advance.
\end{itemize}

\begin{tipbox}[Exam Tip: k-means converges when assignments stop changing]
On 3$\times$3 SAD examples, the assignment grid stabilises after one or two iterations. State explicitly: ``Allocation is unchanged, so cluster centres are also unchanged and the algorithm has converged.''
\end{tipbox}

\separator

\begin{defbox}[Agglomerative Hierarchical Clustering]
Bottom-up: each data point starts as its own cluster; iteratively merge the two most similar clusters until a stopping criterion (single cluster, target $k$, or proximity threshold).
\end{defbox}

\textbf{Pseudo-code (agglomerative):}
\begin{enumerate}
  \item Assign each data point to a unique cluster.
  \item Compute the proximity matrix (distance between every pair of clusters).
  \item Repeat:
  \begin{enumerate}
    \item Merge the two closest clusters.
    \item Update the proximity matrix.
  \end{enumerate}
  \item Until only a single cluster remains (or target reached).
\end{enumerate}

\textbf{Inter-cluster similarity / linkage methods:}

\begin{center}
\renewcommand{\arraystretch}{1.3}
\begin{tabular}{@{} l p{8.5cm} @{}}
\toprule
\textbf{Linkage} & \textbf{Distance between clusters} \\
\midrule
Single-link & MIN distance between any pair of elements \\
Complete-link & MAX distance between any pair of elements \\
Group-average & AVERAGE distance over all pairs of elements \\
Centroid & Distance between the cluster centroids (means) \\
\bottomrule
\end{tabular}
\end{center}

Different linkages produce different dendrograms on the same data.

\begin{tipbox}[Exam Tip: ties merge together]
In tutorial questions on agglomerative clustering, several pairs may be tied at the minimum distance. The model answer merges \emph{all tied pairs in the same step}, then updates the proximity matrix once.
\end{tipbox}

\separator

\begin{defbox}[Segmentation by Graph Cutting]
Build a weighted graph $G = (V, E)$. Each vertex $v \in V$ is an image element (a feature vector); each edge $e_{pq} \in E$ has weight equal to the similarity between connected elements. Segmentation $=$ partitioning $G$ into disjoint subgraphs.
\end{defbox}

For two disjoint sets $A, B \subseteq V$:
\[
\mathrm{cut}(A, B) = \sum_{p \in A,\, q \in B} e_{pq}
\]
Cutting edges with minimum total weight separates the most dissimilar parts of the graph. \textbf{Bias problem:} minimum cut favours small subgraphs (cost grows with the number of cut edges).

\begin{thmbox}[Normalised Cuts (Shi \& Malik)]
Normalise the cut cost by the total association of each side with the whole graph:
\[
\mathrm{Ncut}(A, B) = \frac{\mathrm{cut}(A, B)}{\mathrm{assoc}(A, V)} + \frac{\mathrm{cut}(A, B)}{\mathrm{assoc}(B, V)},
\quad
\mathrm{assoc}(A, V) = \sum_{p \in A,\, q \in V} e_{pq}
\]
Minimising Ncut avoids the bias toward small subgraphs.
\end{thmbox}

\textbf{Pseudo-code (Normalised Cuts):}
\begin{enumerate}
  \item Form a graph in which each vertex is an image element and each edge weight is the similarity between connected elements.
  \item For every possible bipartition $(A, B)$ of vertices, compute $\mathrm{Ncut}(A, B)$.
  \item Cut the edges between $A$ and $B$ for which $\mathrm{Ncut}$ is minimum.
  \item Repeat steps 2--3 for each subgraph until either $\mathrm{Ncut}$ exceeds a threshold for every subgraph, or the maximum number of segments is reached.
\end{enumerate}

\textbf{Problems with Ncuts:}
\begin{itemize}
  \item Finding the exact minimum is NP-hard --- in practice we use approximate (eigenvector-based) solutions.
  \item Bias toward partitioning into roughly equal segments.
  \item Struggles with textured backgrounds (a common weakness across most segmentation algorithms).
\end{itemize}

\bigseparator

% -----------------------------------------------------------------------------
\subsubsection{Fitting Methods}

\begin{defbox}[Fitting]
A class of segmentation methods that fit a mathematical model (line, circle, arbitrary parametric shape, smooth curve) to image elements. Elements that fit the model are grouped together; the model's boundary becomes the segmentation.
\end{defbox}

\separator

\begin{defbox}[Hough Transform (for straight lines)]
For each candidate edge pixel, compute the parameters $(\theta, r)$ of \emph{every} line passing through that pixel and increment a vote in an accumulator array indexed by $(\theta, r)$. Peaks in the accumulator correspond to lines supported by many pixels.
\end{defbox}

\textbf{Line parameterisation.} Avoid $y = mx + c$ (slope $m$ and intercept $c$ go to infinity for vertical lines). Use angle $\theta$ (with the horizontal axis) and perpendicular distance $r$ from the origin:
\[
y \cos\theta - x \sin\theta = r
\]
Note: in image coordinates the $y$-axis points downwards, so positive $\theta$ corresponds to a clockwise rotation.

\textbf{Pseudo-code (Hough transform for straight lines):}
\begin{enumerate}
  \item Initialise all accumulator entries to 0.
  \item For each edge pixel $(x, y)$:
  \begin{enumerate}
    \item For each quantised value of $\theta$:
    \begin{enumerate}
      \item Compute $r = y \cos\theta - x \sin\theta$.
      \item Increment the accumulator entry nearest $(\theta, r)$.
    \end{enumerate}
  \end{enumerate}
  \item Find peaks in the accumulator: each peak is a line.
\end{enumerate}

\textbf{Sizing the accumulator.} The maximum $|r|$ is the image diagonal $\sqrt{W^2 + H^2}$ (round up). Quantise $\theta$ and $r$ as required by the question.

\separator

\textbf{Generalised Hough Transform.}
\begin{itemize}
  \item \textbf{Circles}: parameterise by centre $(a, b)$ and radius $r$ --- 3-D accumulator. Each edge pixel votes for all $(a, b, r)$ corresponding to circles passing through it.
  \item \textbf{Arbitrary (non-parametric) shapes}: store a table of edge offsets relative to a reference point. Each image edge pixel assumes it could be each table entry in turn and votes for the implied reference-point location. Peaks in the reference-point accumulator identify shape locations.
\end{itemize}

\separator

\begin{thmbox}[Hough Transform: Trade-offs]
\textbf{Advantages:}
\begin{itemize}
  \item Tolerant to gaps in edges and to occlusion.
  \item Relatively unaffected by noise.
  \item Detects multiple instances of the shape in a single pass.
\end{itemize}
\textbf{Disadvantages:}
\begin{itemize}
  \item Expensive in memory and computation (especially for $\geq$ 3 parameters).
  \item Localisation information lost (line endpoints unknown).
  \item Peaks hard to locate with noisy or parallel edges.
  \item Sensitive to accumulator quantisation: too coarse $\Rightarrow$ different lines merge; too fine $\Rightarrow$ noise scatters votes and lines are missed.
\end{itemize}
\end{thmbox}

\begin{tipbox}[Exam Tip: quantisation matters]
Tutorial Q19 reuses the same image as Q18 but with $\theta \in \{0, 30, 60, 90, 120, 150\}$ instead of $\{0, 45, 90, 135\}$, and there is no clear peak any more. Be ready to comment that ``results are sensitive to the quantisation of the parameters''.
\end{tipbox}

\separator

\begin{defbox}[Active Contours (``Snakes'')]
A snake is a spline curve that deforms to minimise an energy function combining a shape (internal) term and an image (external) term. The minimum-energy curve is short, smooth, and lies along intensity discontinuities.
\end{defbox}

\textbf{Energy function:}
\[
E_{\text{snake}} = E_{\text{internal}} + E_{\text{external}}
\]
\begin{itemize}
  \item \textbf{Internal energy}: function of the snake's \emph{shape}. Reduced when the curve is short and smooth (penalises bending and stretching).
  \item \textbf{External energy}: function of \emph{image features} near the contour. Reduced when intensity gradients along the curve are high (i.e.\ the curve sits on edges).
\end{itemize}

\textbf{Behaviour.} Initialise around an object; the snake ``shrink-wraps'' onto the boundary as energy is minimised.

\textbf{Problems:}
\begin{itemize}
  \item Only works for shapes that are smooth and form a closed contour.
  \item Extremely sensitive to parameters that weight internal vs.\ external energy.
\end{itemize}

\bigseparator

% =============================================================================
\subsection{Worked Examples}

\begin{exbox}[Example 1: Role of Mid-Level Vision (Q1)]
\textbf{Q:} Briefly describe the role of mid-level vision.

\textbf{A:} To group together image elements that belong together, and to segment them from all other image elements.
\end{exbox}

\begin{exbox}[Example 2: Hysteresis Thresholding (Q3)]
\textbf{Q:} Apply hysteresis thresholding to the 3$\times$3 image $I$ with $t_h = 0.75$, $t_l = 0.25$, 4-connected neighbours.
\[
I = \begin{pmatrix} 0.4 & 0.1 & 0.4 \\ 0.2 & 0.2 & 0.7 \\ 0.3 & 0.8 & 0.9 \end{pmatrix}
\]

\textbf{A:} Step 1 --- pixels with $I > t_h$ become 1, pixels with $I < t_l$ become 0; mid-range pixels (those with $t_l \leq I \leq t_h$, including $0.4, 0.7, 0.3$) are left undecided:
\[
\begin{pmatrix} 0.4 & 0 & 0.4 \\ 0 & 0 & 0.7 \\ 0.3 & 1 & 1 \end{pmatrix}
\]
Step 2 --- a mid-range pixel becomes 1 \emph{only if it neighbours a pixel already classified as figure} (4-connected). The seeds (1s) are at $(3,2)$ and $(3,3)$. Propagating:
\begin{itemize}
  \item $0.7$ at $(2,3)$: 4-neighbour $(3,3)=1$ $\Rightarrow$ becomes 1.
  \item $0.3$ at $(3,1)$: 4-neighbour $(3,2)=1$ $\Rightarrow$ becomes 1.
  \item $0.4$ at $(1,1)$: 4-neighbours $\{0.1, 0.2\}$ are all 0 $\Rightarrow$ becomes 0.
  \item $0.4$ at $(1,3)$: 4-neighbours $\{0.1, 0.7\}$ --- the $0.7$ neighbour was just promoted to 1, but the model answer treats Step~2 as a single pass over the original mid-range set, so $0.4$ at $(1,3)$ stays 0 (neighbour $0.1=0$, $0.7$ not yet finalised when row 1 is checked).
\end{itemize}
\textbf{Final binary image:}
\[
I' = \begin{pmatrix} 0 & 0 & 0 \\ 0 & 0 & 1 \\ 1 & 1 & 1 \end{pmatrix}
\]
\end{exbox}

\begin{exbox}[Example 3: Opening vs Closing (Q4)]
\textbf{Q:} For a 7$\times$9 binary image: (a) draw the result of erosion then dilation with 4-connected neighbours; (b) draw the result of dilation then erosion with 8-connected neighbours.

\textbf{A:}
\begin{itemize}
  \item \textbf{(a) Opening (4-conn)}: erosion first removes thin isolated foreground pixels and bridges (any FG pixel touching a BG pixel disappears). Dilation then re-grows the surviving cores. Net effect: small specks of foreground noise are removed; thicker shapes are preserved with roughly their original silhouette.
  \item \textbf{(b) Closing (8-conn)}: dilation first fills small holes and gaps (every BG pixel adjacent to FG --- including diagonals --- becomes FG). Erosion then trims the outer halo back. Net effect: foreground holes/cracks are filled; outer shape is preserved.
\end{itemize}

\begin{tipbox}[Exam Tip: When asked to draw, sketch the foreground silhouette]
Mark which pixels switch on each pass. With 8-connectivity, dilation reaches diagonally too; opening with 4-connectivity is more aggressive than with 8-connectivity (more pixels touch the background under 4-conn).
\end{tipbox}
\end{exbox}

\begin{exbox}[Example 4: Region Growing on a 3$\times$3 Image (Q6)]
\textbf{Q:} Apply region growing with SAD-similarity threshold $T = 12$, 8-connected neighbours, seed $=$ top-left, to the 3$\times$3 feature-vector image:
\[
\begin{matrix} (5,10,15) & (10,15,30) & (10,10,25) \\ (10,10,15) & (5,20,15) & (10,5,30) \\ (5,5,15) & (30,10,5) & (30,10,10) \end{matrix}
\]

\textbf{A:} Use SAD between feature vectors. Starting from the top-left (label 1):
\begin{itemize}
  \item SAD$($top-left, mid-left$) = 5$, SAD$($top-left, top-mid$) = 25$, SAD$($top-left, mid-mid$) = 10$. Add mid-left (5) and mid-mid (10) to region 1.
  \item Grow from mid-left: SAD$($mid-left, bot-left$) = 10$ $\Rightarrow$ add bot-left to region 1.
  \item No further growth from region 1 (all remaining neighbours have SAD $\geq 25$). Pick a new seed (top-mid, label 2).
  \item Grow from top-mid: SAD$($top-mid, top-right$) = 10$ $\Rightarrow$ add. From top-right SAD$($top-right, mid-right$) = 10$ $\Rightarrow$ add. Region 2 cannot grow further.
  \item Pick last seed bot-right (label 3): SAD$($bot-right, bot-mid$) = 5$ $\Rightarrow$ add bot-mid to region 3.
\end{itemize}
\textbf{Final labels:}
\[
\begin{pmatrix} 1 & 2 & 2 \\ 1 & 1 & 2 \\ 1 & 3 & 3 \end{pmatrix}
\]
\end{exbox}

\begin{exbox}[Example 5: k-means on a 3$\times$3 Image (Q12)]
\textbf{Q:} Apply k-means with SAD distance, $k=2$, initial centres $c_1 = (5,10,15)$, $c_2 = (10,10,25)$ to the same 3$\times$3 feature-vector image as Example 4.

\textbf{A:} \textbf{Iteration 1.} Compute SAD distance from each pixel to each centre:
\begin{itemize}
  \item Distances to $c_1 = (5,10,15)$: $\begin{pmatrix} 0 & 25 & 15 \\ 5 & 10 & 25 \\ 5 & 35 & 30 \end{pmatrix}$
  \item Distances to $c_2 = (10,10,25)$: $\begin{pmatrix} 15 & 10 & 0 \\ 10 & 25 & 10 \\ 20 & 40 & 35 \end{pmatrix}$
\end{itemize}
Allocate to closest centre:
$\begin{pmatrix} 1 & 2 & 2 \\ 1 & 1 & 2 \\ 1 & 1 & 1 \end{pmatrix}$.

\textbf{Recompute centres:}
\[
c_1 = \tfrac{1}{6}\big[(5,10,15)+(10,10,15)+(5,20,15)+(5,5,15)+(30,10,5)+(30,10,10)\big] = (14.17, 10.83, 12.50)
\]
\[
c_2 = \tfrac{1}{3}\big[(10,15,30)+(10,10,25)+(10,5,30)\big] = (10, 10, 28.33)
\]

\textbf{Iteration 2.} Recompute distances; allocations are unchanged $\Rightarrow$ centres unchanged $\Rightarrow$ algorithm has converged.
\end{exbox}

\begin{exbox}[Example 6: Linkage Methods in Agglomerative Clustering (Q14)]
\textbf{Q:} Briefly describe (a) single-link, (b) complete-link, (c) group-average, (d) centroid linkage.

\textbf{A:}
\begin{itemize}
  \item \textbf{Single-link}: distance between clusters is the shortest distance between any pair of elements (MIN).
  \item \textbf{Complete-link}: longest distance between any pair of elements (MAX).
  \item \textbf{Group-average}: average distance over all pairs of elements (AVERAGE).
  \item \textbf{Centroid}: distance between the cluster centroids (means of feature vectors).
\end{itemize}
Different linkages can produce dramatically different dendrograms on the same data.
\end{exbox}

\begin{exbox}[Example 7: Hough Transform for a 4$\times$3 Binary Image (Q18)]
\textbf{Q:} The image has foreground pixels at $(0,0), (1,1), (2,2), (3,0)$ (using top-left origin, $y$ down). Compute the Hough accumulator for $\theta \in \{0, 45, 90, 135\}$ deg, $r$ rounded to nearest integer.

\textbf{A:} Recall $r = y\cos\theta - x\sin\theta$. Maximum $|r| \leq \sqrt{3^2 + 2^2} \approx 3.6$, so $r \in [-4, 4]$.

For each pixel, compute the four $(\theta, r)$ pairs:
\begin{center}
\renewcommand{\arraystretch}{1.2}
\begin{tabular}{@{} l c c c c @{}}
\toprule
\textbf{Pixel} & $\theta=0$ & $\theta=45$ & $\theta=90$ & $\theta=135$ \\
\midrule
$(0,0)$ & $r=0$ & $r=0$ & $r=0$ & $r=0$ \\
$(1,1)$ & $r=1$ & $r=0$ & $r=-1$ & $r\approx -1$ \\
$(2,2)$ & $r=2$ & $r=0$ & $r=-2$ & $r\approx -3$ \\
$(3,0)$ & $r=0$ & $r\approx -2$ & $r=-3$ & $r\approx -2$ \\
\bottomrule
\end{tabular}
\end{center}

The maximum accumulator entry is $3$ at $(\theta, r) = (45, 0)$, contributed by the three diagonal pixels $(0,0), (1,1), (2,2)$. This is the line of three collinear edge pixels.

\textbf{Follow-up (Q19).} Repeat with $\theta \in \{0, 30, 60, 90, 120, 150\}$: \emph{no clear peak emerges}, illustrating the algorithm's sensitivity to angle quantisation.
\end{exbox}

\begin{exbox}[Example 8: Active Contour Energy (Q20)]
\textbf{Q:} Briefly describe the two terms that form the energy function that an active contour attempts to minimise.

\textbf{A:} $E_{\text{snake}} = E_{\text{internal}} + E_{\text{external}}$.
\begin{itemize}
  \item \textbf{Internal energy} is a function of the shape of the contour; it is reduced if the curve is short and smooth (i.e.\ penalises bending and stretching).
  \item \textbf{External energy} is a function of the image features near the contour; it is reduced if the intensity gradient along the contour is high (i.e.\ the curve lies on strong edges).
\end{itemize}
\end{exbox}

\bigseparator

% =============================================================================
\subsection{Summary}

\begin{keybox}[Key Takeaways]
\begin{enumerate}
  \item \textbf{Mid-level vision = grouping}. Choose a feature, define similarity (SAD/SSD/Euclidean) or affinity, partition feature space.
  \item \textbf{Thresholding} segments a 1-D feature space. Use the histogram valley; for noisy boundaries use \term{hysteresis} with two thresholds; clean up with morphology (open removes specks, close fills holes).
  \item \textbf{Region-based} methods use spatial information. Region growing (seed-driven), region merging (bottom-up, order-dependent), split-and-merge (top-down then bottom-up via a quadtree).
  \item \textbf{k-means} is partitional clustering: needs $k$, sensitive to initialisation, struggles with non-globular or unequal-size clusters.
  \item \textbf{Agglomerative hierarchical clustering} merges nearest clusters; linkage choice (single/complete/average/centroid) shapes the dendrogram.
  \item \textbf{Graph cuts / Normalised Cuts} treat segmentation as graph partitioning. Min-cut is biased to small subgraphs; Ncut normalises by association. Optimum is NP-hard.
  \item \textbf{Hough transform} fits parametric shapes by voting; tolerant to gaps and occlusion but expensive and sensitive to quantisation. Use $r = y\cos\theta - x\sin\theta$.
  \item \textbf{Snakes} minimise internal (smoothness) plus external (image-gradient) energy; only suitable for smooth closed boundaries.
\end{enumerate}
\end{keybox}

\textbf{Comparison of segmentation methods:}

\begin{center}
\renewcommand{\arraystretch}{1.3}
\begin{tabular}{@{} l p{3.6cm} p{3.6cm} p{3.6cm} @{}}
\toprule
\textbf{Method} & \textbf{Feature Space} & \textbf{Spatial Info?} & \textbf{Main Pitfall} \\
\midrule
Thresholding & 1-D (intensity / filter output) & No & Single threshold rarely works globally \\
Hysteresis thresholding & 1-D + adjacency & Yes (adjacency) & Choice of $t_l, t_h$ \\
Region growing & multi-D + adjacency & Yes & Leaks through weak boundaries \\
Region merging & multi-D + adjacency & Yes & Order-dependent results \\
Split-and-merge & multi-D + quadtree adjacency & Yes & Quadtree artefacts; padding needed \\
k-means & multi-D & Optional (add position) & Init-dependent; needs $k$; non-globular fails \\
Agglomerative & multi-D & Optional & Linkage choice changes results; $O(n^2)$ memory \\
Graph cuts (min-cut) & graph in feature space & Yes (via edges) & Bias toward small subgraphs \\
Normalised Cuts & graph in feature space & Yes & NP-hard exact; equal-size bias; texture issues \\
Hough transform & parameter space (vote) & Yes & Memory cost; quantisation sensitivity \\
Active contours (snakes) & curve over image & Yes (along curve) & Smooth closed shapes only; parameter-sensitive \\
\bottomrule
\end{tabular}
\end{center}

\begin{tipbox}[Exam Tip: Pseudo-code is a recurring marks-grab]
Many tutorial questions ask you to write pseudo-code for an algorithm (region growing, region merging, split-and-merge, k-means, agglomerative, Ncuts, Hough). Memorise the canonical step lists in this chapter.
\end{tipbox}
